% !TeX program = xelatex
% !BIB program = biber
% 格式案例：临时将Main document切换为example.tex即可预览。
% 正常写作时Main document选择main.tex；编译器选择XeLaTeX。
% 本文件使用muban/lunwen.cls，另附示例文献库和两张示意图。
\def\mulu{./}
\makeatletter
\edef\input@path{{\mulu muban/}{\mulu}}
\makeatother
\documentclass{lunwen}
% 以下文献均为虚构的排版演示条目，保存在muban/example-wenxian.bib中，不改正式文献库。

\edef\wenxianlujing{\mulu muban/example-wenxian.bib}
\expandafter\addbibresource\expandafter{\wenxianlujing}

% 分图编号采用全角括号，分图题保持五号字并放在分图下方。
\DeclareCaptionLabelFormat{fenhao}{（#2）}
\captionsetup[subfigure]{font=thesis,labelformat=fenhao,
 labelsep=thesisspace,justification=centering,singlelinecheck=false}

% 缺失值按两个正文数字的宽度居中。
\newlength{\queshikuan}
\newcommand{\queshizhi}{\makebox[\queshikuan][c]{—}}

% 符号说明：正文续行对齐到破折号后的文字。
\newcommand{\shiliangshuoming}[3]{%
 \par\begingroup\hangindent=5.5em\hangafter=1\noindent
 \makebox[2em][l]{#1}\makebox[1.5em][l]{$#2$}%
 \makebox[2em][l]{——}#3\par\endgroup}

\begin{document}
\frontmatter\mainmatter
\section{常用格式示例}

% 标题层级示例：四级以内，正文首句不额外加粗。
\subsection{标题层级与正文}
本节演示中文正文、英文Architecture与数字0123456789混排。
正文首句保持常规字重，中文说明使用全角标点。
\subsubsection{三级标题示例}
三级标题下正常书写正文。
\paragraph{四级标题示例}
四级标题以下不再增加第五级标题。

% 一、图片插入与图片引用。
\subsection{图片插入与引用}
本文件中的图表数据和文献条目均用于排版演示；正式写作时替换为真实内容与来源。

研究流程见\figref{fig:liucheng}。图片编号由模板自动生成。

\begin{figure}[htbp]
  \centering
  % 修改下面的文件路径即可换图；.75表示正文宽度的75%。
  \includegraphics[width=.75\linewidth]{\mulu tupian/example-liucheng.pdf}
  \figurenote{注：此图为示意图，资料来源为作者绘制}
  \caption{研究流程示意}
  \label{fig:liucheng}
\end{figure}

引用图片使用以下命令：
\begin{verbatim}
\figref{fig:liucheng}
\end{verbatim}
图的标签写在\verb|\caption|后面。

% 二、文献引用。使用现有文献库中的两个条目。
\clearpage
\subsection{文献引用}
首次引用期刊文章\cite{case-journal}，随后引用专著\cite{case-book}。
句中引用可以写作：由文献\inlinecite{case-conference}可知。

学位论文、论文集析出文献和网页文献分别引用为
\cite{case-thesis}、\cite{case-chapter}和\cite{case-web}。
网页文献的示例条目同时填写网址与访问日期。

同一文献不同页码的引用为\cite[15-18]{case-journal}，
再次引用其他页码为\cite[22]{case-journal}，文献编号不变。
非连续与连续文献合并引用可写为\cite{case-journal,case-conference,case-thesis,case-chapter}。

\begin{verbatim}
\cite{case-journal}
\inlinecite{case-conference}
\cite[15-18]{case-journal}
\cite{case-journal,case-conference,case-thesis,case-chapter}
\end{verbatim}

参考文献末页分别演示期刊、专著、会议、学位论文、
论文集析出文献和网页的著录。示例数量不代表正式论文的文献数量要求。

% % 图片补充：分图、引用图题及续图。
\clearpage
\subsection{分图及来源标注}
分图案例见\figref{fig:fentu}，其中左侧为
\figref{fig:fentu}（\subref{fig:fentu-a}），右侧为
\figref{fig:fentu}（\subref{fig:fentu-b}）。

\begin{figure}[htbp]
 \centering
 \begin{subfigure}[t]{.46\linewidth}
  \centering
  \includegraphics[width=\linewidth]{\mulu tupian/example-fentu.pdf}
  \figurenote{注：左侧方案说明}
  \caption{方案甲}\label{fig:fentu-a}
 \end{subfigure}\hfill
 \begin{subfigure}[t]{.46\linewidth}
  \centering
  \includegraphics[width=\linewidth]{\mulu tupian/example-fentu.pdf}
  \figurenote{注：右侧方案说明}
  \caption{方案乙}\label{fig:fentu-b}
 \end{subfigure}
 \figurenote{注：分图说明位于各分图题上方，总说明位于总图题上方}
 \caption{分图排版示例}\label{fig:fentu}
\end{figure}


引用来源应先在引出图片的正文中出现\cite{case-journal}，
并在引用图的图题右上角再次标注，见\figref{fig:yuantu}。
本例的图片是示意图，仅演示来源标注位置。

\begin{figure}[htbp]
 \centering
 \includegraphics[width=.7\linewidth]{\mulu tupian/example-liucheng.pdf}
 \figurenote{注：图题来源标注位置演示}
 \caption{引用图题格式示例\cite{case-journal}}\label{fig:yuantu}
\end{figure}


\clearpage
\subsection{跨页分图与续图}
分图较多而需要跨页时，使用同一个图号，
总图题只在第一页出现。示例见\figref{fig:xutu}。

\begin{figure}[htbp]
 \centering
 \begin{subfigure}{.46\linewidth}
  \centering\includegraphics[width=\linewidth]{\mulu tupian/example-fentu.pdf}
  \caption{第一组方案}
 \end{subfigure}\hfill
 \begin{subfigure}{.46\linewidth}
  \centering\includegraphics[width=\linewidth]{\mulu tupian/example-fentu.pdf}
  \caption{第二组方案}
 \end{subfigure}
 \caption{多分图跨页排版示例}\label{fig:xutu}
\end{figure}

\clearpage
% 续图页顶预留间距，避免图形接触页眉双线。
\vspace*{8bp}
\begin{figure}[htbp]
 \ContinuedFloat
 \setcounter{subfigure}{2}
 \centering
 \begin{subfigure}{.46\linewidth}
  \centering\includegraphics[width=\linewidth]{\mulu tupian/example-fentu.pdf}
  \caption{第三组方案}
 \end{subfigure}\hfill
 \begin{subfigure}{.46\linewidth}
  \centering\includegraphics[width=\linewidth]{\mulu tupian/example-fentu.pdf}
  \caption{第四组方案}
 \end{subfigure}
 \captionsetup{labelsep=none,list=no}
 \caption{（续图）}
\end{figure}


\clearpage
\subsection{坐标单位及地图要素}
坐标图需要同时标明数值和相应计量单位，见\figref{fig:zuobiao}。

\begin{figure}[htbp]
 \centering\tablefont
 \setlength{\unitlength}{1bp}
 \begin{picture}(290,160)
  \put(30,25){\vector(1,0){225}}
  \put(30,25){\vector(0,1){110}}
  \put(0,149){\makebox(0,0)[l]{高度（$\mathrm{m}$）}}
  \put(190,5){\makebox(0,0)[l]{距离（$\mathrm{m}$）}}
  \put(23,13){\makebox(0,0){0}}
  \put(105,25){\line(0,1){4}}\put(105,13){\makebox(0,0){50}}
  \put(180,25){\line(0,1){4}}\put(180,13){\makebox(0,0){100}}
  \put(30,75){\line(1,0){4}}\put(18,75){\makebox(0,0){10}}
  \put(30,125){\line(1,0){4}}\put(18,125){\makebox(0,0){20}}
  \put(75,45){\circle*{3}}\put(120,70){\circle*{3}}
  \put(165,95){\circle*{3}}\put(210,120){\circle*{3}}
 \end{picture}
 \caption{坐标与单位示例}\label{fig:zuobiao}
\end{figure}


地图要素的位置示意见\figref{fig:ditu}。
此图不含真实地理底图，仅演示研究区、图例、指北针和比例尺的位置。
正式地图应替换为合规底图，并检查领土完整与底图使用说明。

\begin{figure}[htbp]
 \centering\tablefont\setlength{\unitlength}{1bp}
 \begin{picture}(290,165)
  \put(20,35){\framebox(220,100){}}
  \put(32,118){\makebox(0,0)[l]{示意底图}}
  \put(80,65){\framebox(90,35){研究区}}
  \put(260,80){\vector(0,1){55}}
  \put(260,147){\makebox(0,0){北}}
  \put(35,17){\line(1,0){100}}
  \put(35,14){\line(0,1){6}}
  \put(85,14){\line(0,1){6}}
  \put(135,14){\line(0,1){6}}
  \put(35,5){\makebox(0,0){0}}
  \put(85,5){\makebox(0,0){50}}
  \put(135,5){\makebox(0,0){100 $\mathrm{m}$}}
  \put(190,12){\framebox(10,7){}}
  \put(206,15){\makebox(0,0)[l]{研究区}}
 \end{picture}
 \caption{地图要素位置示意}\label{fig:ditu}
\end{figure}


% 三、普通三线表与表格引用。
\clearpage
\subsection{三线表制作与引用}
研究变量见\tabref{tab:bianliang}。表题位于表格上方。

\begin{table}[htbp]
  \centering\tablefont
  \caption{研究变量说明}
  \label{tab:bianliang}
  \begin{tabularx}{\linewidth}{lXr}
    \toprule
    变量 & 含义 & 示例值\\
    \midrule
    Area & 研究单元面积（$\mathrm{m}^2$） & 1200\\
    Count & 观测数量 & 36\\
    Score & 综合得分 & 3.14\\
    Other & 缺失值 & —\\
    \bottomrule
  \end{tabularx}
\end{table}

引用表格使用以下命令：
\begin{verbatim}
\tabref{tab:bianliang}
\end{verbatim}
表格标签写在\verb|\caption|后面。

三条横线分别由\verb|\toprule|、\verb|\midrule|、\verb|\bottomrule|生成。
单元格之间用\verb|&|分隔，每行以\verb|\\|结束。

% 表格补充：统一单位、合并单元格和引用表题。
\clearpage
\subsection{统一单位与合并单元格}
全表数值均为面积时，将单位放在表头右上方，见\tabref{tab:danwei}。
表内缺失值按两个数字的宽度居中。

\begin{table}[htbp]
 \centering\tablefont
 \settowidth{\queshikuan}{00}
 \caption{统一单位与单元格合并示例}\label{tab:danwei}
 \begin{tabularx}{\linewidth}{lXrr}
  \multicolumn{4}{r}{（$\mathrm{m}^2$）}\\
  \toprule 分组 & 编号 & 更新前 & 更新后\\\midrule
  \multirow{2}{*}{甲组} & A01 & 1200 & 1500\\
   & A02 & 1300 & 1600\\
  乙组 & B01 & \multicolumn{2}{c}{\queshizhi}\\
  \multicolumn{2}{l}{合计} & 2500 & 3100\\
  \bottomrule
 \end{tabularx}
\end{table}

纵向合并使用\verb|\multirow|，横向合并使用\verb|\multicolumn|。
本表的分组纵向合并，合计标签及缺失区域横向合并。

引用文献中的表格，应先在正文标注来源\cite{case-journal}，
再在表题右上角标注，见\tabref{tab:yinyong}。
Word硕士模板通常只用中文表题；本例额外显示英文行，
用于演示PPT提及的英文表题来源标注位置，普通表格保持中文表题。

\begin{table}[htbp]
 \centering\tablefont
 \caption{引用表题格式示例\cite{case-journal}}\label{tab:yinyong}
 \caption*{Table \thetable\quad Example of source citation\cite{case-journal}}
 \begin{tabularx}{\linewidth}{Xr}
  \toprule 对象 & 示意数值\\\midrule
  测试甲 & 10\\
  测试乙 & 20\\\bottomrule
 \end{tabularx}
\end{table}


% 四、跨页表格。续表沿用“表号（续表）”右对齐的形式。
\clearpage
\subsection{跨页表格与续表}
较长的数据表使用\verb|longtable|，见\tabref{tab:changbiao}。
本例通过足够多的数据行自然产生续页，续表标题保持右对齐。
\begingroup
\tablefont
\begin{longtable}{p{25mm}p{108mm}}
  % 第一页的表题及表头。
  \caption{研究单元记录}\label{tab:changbiao}\\
  \toprule
  研究单元 & 观测记录\\
  \midrule
  \endfirsthead

  % 续页的表题及表头。8bp只调整续页表头的上方留白。
  \noalign{\vskip 8bp}
  \multicolumn{2}{r}{表\thetable（续表）}\\
  \toprule
  研究单元 & 观测记录\\
  \midrule
  \endhead

  % 每页表格底线。
  \bottomrule
  \endfoot

  单元01 & 开放空间面积为1200平方米\\
  单元02 & 开放空间面积为1250平方米\\
  单元03 & 开放空间面积为1300平方米\\
  单元04 & 开放空间面积为1350平方米\\
  单元05 & 开放空间面积为1400平方米\\
  单元06 & 开放空间面积为1450平方米\\
  单元07 & 开放空间面积为1500平方米\\
  单元08 & 开放空间面积为1550平方米\\
  单元09 & 开放空间面积为1600平方米\\
  单元10 & 开放空间面积为1650平方米\\
  单元11 & 开放空间面积为1700平方米\\
  单元12 & 开放空间面积为1750平方米\\
  单元13 & 开放空间面积为1800平方米\\
  单元14 & 开放空间面积为1850平方米\\
  单元15 & 开放空间面积为1900平方米\\
  单元16 & 开放空间面积为1950平方米\\
  单元17 & 开放空间面积为2000平方米\\
  单元18 & 开放空间面积为2050平方米\\
  单元19 & 开放空间面积为2100平方米\\
  单元20 & 开放空间面积为2150平方米\\
  单元21 & 开放空间面积为2200平方米\\
  单元22 & 开放空间面积为2250平方米\\
  单元23 & 开放空间面积为2300平方米\\
  单元24 & 开放空间面积为2350平方米\\
  单元25 & 开放空间面积为2400平方米\\
  单元26 & 开放空间面积为2450平方米\\
  单元27 & 开放空间面积为2500平方米\\
  单元28 & 开放空间面积为2550平方米\\
  单元29 & 开放空间面积为2600平方米\\
  单元30 & 开放空间面积为2650平方米\\
  单元31 & 开放空间面积为2700平方米\\
  单元32 & 开放空间面积为2750平方米\\
  单元33 & 开放空间面积为2800平方米\\
  单元34 & 开放空间面积为2850平方米\\
  单元35 & 开放空间面积为2900平方米\\
  单元36 & 开放空间面积为2950平方米\\
  单元37 & 开放空间面积为3000平方米\\
  单元38 & 开放空间面积为3050平方米\\
  单元39 & 开放空间面积为3100平方米\\
  单元40 & 开放空间面积为3150平方米\\
\end{longtable}
\endgroup

% 五、公式及公式引用。
\clearpage
\subsection{公式制作与引用}
行内公式直接放在正文中，例如$P_i=N_i/A_i$。

独立编号公式见\eqnref{eq:zhibiao}。
\begin{equation}
  P_i=\frac{N_i}{A_i}\qquad i=1,2,\ldots,10
  \label{eq:zhibiao}
\end{equation}
\shiliangshuoming{式中}{P_i}{第$i$个研究单元的指标值。这里使用较长的说明文字演示自动换行，续行应对齐到破折号后第一个汉字的位置，而不是回到整段的左边界。}
\shiliangshuoming{}{N_i}{第$i$个研究单元的观测数量。}
\shiliangshuoming{}{A_i}{第$i$个研究单元的面积，单位为$\mathrm{m}^2$。}
公式引用使用\verb|\eqnref{eq:zhibiao}|，标签写在公式环境内。

下式同时演示求和、根号、上下标及主副分数线：
\begin{equation}
  Q_i=\frac{\dfrac{\sum_{j=1}^{n}w_jx_{ij}}{\sum_{j=1}^{n}w_j}}{\sqrt{1+\alpha^2}}
  \label{eq:zonghe}
\end{equation}
\shiliangshuoming{式中}{Q_i}{第$i$个研究单元的综合指标。}
\shiliangshuoming{}{w_j}{第$j$项指标的权重，权重总和不为零。}
\shiliangshuoming{}{x_{ij}}{第$i$个研究单元的第$j$项指标观测值。}
\shiliangshuoming{}{n}{参与计算的指标总数。}
\shiliangshuoming{}{\alpha}{调节参数。}

多行公式共用一个编号时，使用\verb|aligned|：
\begin{equation}
  \begin{aligned}
    S_i &= \frac{\sum_{j=1}^{n}w_jx_{ij}}{\sum_{j=1}^{n}w_j}\\
    Q_i &= \frac{S_i}{\sqrt{1+\alpha^2}}
  \end{aligned}
  \label{eq:duohang}
\end{equation}
\shiliangshuoming{式中}{S_i}{第$i$个研究单元的加权平均值，其余符号含义同前。}
再次引用时写作\eqnref{eq:duohang}。

\subsection{本章小结}
本章展示了正文中常用的图表、引用和公式排版方式。小结保持纯文字归纳。

\printbibliography[heading=bibintoc,title={参考文献}]
\end{document}
